The current manuscript has several important limitations. First, the artifact is a prototype baseline with canonical documentation, machine-readable contracts, fixtures, and evaluation scaffolding. It is not yet a completed large-scale empirical study over real production traffic. As a result, some claims in this paper are architectural and methodological rather than conclusively empirical.

Second, early embodied benchmarks may rely on curated or synthetic scenarios before broader field deployment data exists. Those scenarios are still useful for exposing structural planner failures, but they do not automatically capture every complexity of real logistics, local labor coordination, or jurisdiction-specific constraints. Future work should test the model with richer scenario sets and, where feasible, real deployment traces or controlled human studies.

Third, the live-model benchmark introduces another layer of measurement fragility. Provider access limits, output truncation, and JSON-format drift can reduce parse success even when the underlying planning signal is stronger than the exact contract score suggests. We therefore separate call success, parse success, exact pass, and semantic coverage, but the current study is still limited to one gateway environment and one family of model providers.

Fourth, explicit embodied verification introduces ethical and operational concerns. Proof mechanisms such as geostamps, photos, access logs, and witness records can improve accountability, but they can also create privacy, surveillance, or labor-power asymmetries if handled poorly. A request-rooted system should therefore make proof requirements legible and proportional rather than maximizing visibility by default.

Fifth, commercial and legal constraints remain important. Physical work can involve insurance, site safety, labor classification, access rights, or chain-of-custody requirements that are not solved solely by better planning. The object model can preserve those constraints more honestly, but it does not replace the need for policy, compliance, and human judgment.
